\scanpage{45}
(\unsure{ici|i.e.} on peut prendre $\unsure{I|J}=\operatorname{Ouv}(X)$).

\deleted{\unsure{a)}} \uline{NB} Si l'ens. des ouverts de $E$ est \deleted{\unsure{un}} isomorphe (pour la structure ordonnée) à l'ens des ouverts $\operatorname{Ouv}(X)$ quelque $X$, alors $E\to\deleted{\unsure{X}}$ \inserted{(i.e. $E\to\operatorname{Top}(X)$)} \unsure{est|c'est|et} \unsure{l'un|un|c'est} universel pour les morphismes de $E$ dans un espace topologique.

\deleted{c) $\operatorname{Top}(X,G)$ pour \unsure{$(X,G)$ variable} les Morphismes de \inserted{$E\xrightarrow[\substack{\varinjlim\\\varprojlim\ \mathrm{fin.}}]{}\underline{\mathrm{Hom}}(\unsure{\mathrm{Cl}^\circ|C^\circ|\mathcal A(X)},\mathrm{Ens})\Longleftrightarrow\unsure{\mathrm{Cl}^\circ|C^\circ}\to\underline{\mathrm{Hom}}(E,(\mathrm{Ens}))$} \unsure{$\underline{\mathrm{Hom}}(\mathrm{Pt}(E))\Leftrightarrow U^\circ\to\underline{\mathrm{Hom}}(\mathrm{Cl},(\mathrm{Ens}))$}}

c) Morphismes d'un pt dans Topos $E$ : Catégories $\underline{\mathrm{Pt}}(E)\simeq\underline{\mathrm{Fib}}(E)^\circ$

d) \uline{Catégories} $\widehat C$ : Tout \deleted{morphisme} \inserted{foncteur} $u:C\to\unsure{C'|\widehat C'}$ définit un foncteur \unsure{?} en sens inverse \inserted{$u^*_{\deleted{\unsure{*}}}\deleted{\unsure{(F)}}:$} $F'\mapsto F'\circ u$

\[\widehat C'\xrightarrow{u^*}\widehat C\]

\deleted{qui a deux adjoints, notés $u_{\unsure{?|!}}^{\deleted{\unsure{*}}}$ et $u^!$} qui commute \deleted{$u^!\ \unsure{u_*u^*}$} \inserted{$u^*\ u_*\ u_!$} aux $\varprojlim$ quelc. et \deleted{le foncteur $u^*$} $\varinjlim$ quelc., donc est le foncteur image inverse pour un $u:\widehat C\longrightarrow\widehat C'$. \inserted{$u_!\quad u^*\quad u_*$}

% slanted
\margin{donc $\widehat C$ est \uline{covariant} en $C$}

En plus de l'adjoint à droite $u_*$, il a un adjoint à g, qu'on note $u_!$ ....

\uline{Ex} : $B_G$ avec $G$ variable, qui est aussi exemple de 4.1. a)

\unsure{$\leftarrow$|e)} [ Morphismes \inserted{de topos} $\widehat C\longrightarrow\unsure{E|\widehat E}$

\[E\xrightarrow[\substack{\varinjlim\\\varprojlim\ \mathrm{finie}}]{}\underline{\mathrm{Hom}}(C^\circ,\mathrm{Ens})\quad\Longleftrightarrow\quad C^\circ\longrightarrow\underbrace{\underline{\mathrm{Hom}}_{\substack{\varinjlim\\\varprojlim\ \mathrm{finie}}}(E,(\mathrm{Ens}))}_{\parallel\,\underline{\mathrm{Pt}}(E)^\circ}\]

\[\underline{\mathrm{Hom\,top}}(\widehat C,E)\simeq\deleted{\unsure{\underline{\mathrm{Hom}}(C^\circ,\underline{\mathrm{Pt}}(E)^\circ)^\circ|\underline{\mathrm{Hom}}(C^\circ,\underline{\mathrm{Pt}}(E)^\circ)}}\ \underline{\mathrm{Hom}}(C,\underline{\mathrm{Pt}}(E))\]

si $E=\widehat C'$, on trouve $\underline{\mathrm{Hom\,top}}(\widehat C,\widehat C')\simeq\underline{\mathrm{Hom}}(C,\deleted{Pt}\operatorname{Pro}\,C')$ \inserted{$(\underline{\mathrm{Hom}}(C,C')\ \mathrm{pl.fid})$}

[ Si $C$ stable par $\varprojlim$ finies, on trouve

\[\underline{\mathrm{Hom\,top}}(E,\widehat C)\simeq\underline{\mathrm{Hom}}_{\mathrm{ex.g.}}(C,E)^\circ\ \Bigr]\vphantom{\widehat C}\]
% archive sheet number: 44
